\chapter{Introduction}
\label{chap::intro-1}

Athletes in professional, highly competitive sports have every interest in resorting to the newest advances in technology in order to enhance their competitive performance, and training regime. With recent jumps in technology, \gls{vr} is increasingly establishing itself as a valuable tool to professional athletes in training scenarios, thanks to its ability to display a wider variety of settings and activity, at a higher level of control than training grounds ever could grant on a regular basis.\\

Earliest adoptions of VR, or immersive technologies in professional sport include 3D recreations of matches for more immersive match reviews from either first-person, or 3-dimensional non-fixed viewpoints in the early 2010s \cite{farley-2020-virtual}.\\
The 2022 inews article "Half of Premier League clubs use virtual reality to heal injuries and recreate matchday pressure in training"\cite{cunningham-2022-half} highlights the increasing popularity of \gls{vr} across the elite-level football, or soccer, scene, with half of the twenty English Premier League teams  of the time applying \gls{vr} scenarios for injury rehabilitation, recreation of matchday atmosphere and pressures, and other training and academy uses, which were left more vague.\\
The advantages of this technology are manifold: besides the proven ability to positively affect performance and motivation through various different factors interacting together, e.g. immersion (\cite{ijsselsteijn-2004-fun}), competitiveness (\cite{nunes-2014-motivating}), digital cooperation (\cite{irwin-2012-aerobic}), agency in selecting and adjusting the environment (\cite{neumann-2018-systematic}), or adaptive difficulty levels (\cite{gray-2017-transfer}, \cite{richlan-2023-virtual}), it is a technology that allows for safer training conditions, examplified with \cite{ritter-2023-gymnastic} showcasing that gymnasts practising on a beam at floor level - determined to be a safer training environment for the study than the physical facility featuring a beam raised above the floor - whilst wearing a \gls{hmd} experienced similar psychological conditions to the control group practising in the physical environment. \gls{vr} is today also applicable at a small scale, e.g. in local gyms, or in dedicated spaces at home or in training centres, and affordably through commercial-grade \glspl{hmd} \cite{farley-2020-virtual}.\\

With the increased interest, there is a rise in scrutiny of which ways the incorporation of \gls{vr} can benefit athletes, with potential for use in sharpening motor skills such as movement and coordination, cognitive skills like perception on the pitch or positioning, mental skills to better resist match or score pressure, or tactical and strategic proficiency in helping to understand and apply directives from the coach \cite{richlan-2023-virtual}.\\
Newer research, such as \cite{ritter-2023-gymnastic} and \cite{gerwann-2025-immersive}, also promotes the introduction of \gls{av}, which introduces interaction with physical objects to provide or simulate haptic feedback whilst performing tasks in \gls{vr}.\\

However, despite the widening adoption of the technology and increased research, the factors affecting performance in, and skill acquisition through \gls{vr} are still a very young field of research that is yet to accurately study the aspects in a more granular and standardised fashion, as various reviews of existing studies into the topic of performance in \gls{vr} sports applications, and their effect on competitive output, highlight \cite{neumann-2018-systematic}\cite{richlan-2023-virtual}\cite{gerwann-2025-immersive}.\\

Despite a notable extension of the range of sports and sport tasks studied between the reviews of \cite{neumann-2018-systematic} and \cite{richlan-2023-virtual}, the former of which advocated for exactly this extension beyond mostly aerobic workout activities, a sport that seems conspicuously absent from studies so far is Volleyball, alongside any of its variations, and probably for good reasons, as it requires dealing with many inherent limitations of VR application to sport, outlined in \cite{richlan-2023-virtual}.\\

That being said, there is ample evidence that the acquisition and refinement of complex skills in sport does not require a perfect recreation of the environment, and can be compartmentalised to fit specific purposes at a time, as evidenced by \cite{farley-2020-virtual}'s proposal of a system making use of a skateboard with \gls{vr} environments, to practise surfing in dry conditions for example. Therefore, volleyball seems a suitable candidate for further investigation.\\

This dissertation therefore aims to stimulate further research into applying \gls{vr} to volleyball applications by investigating whether, and in what way, simulated crowds meaningfully affect, both positively or negatively, users' sense of immersion or, more specifically, the illusion of plausibility that can make or break this feeling of immersion - itself a primordial factor in promoting the plethora of positive effects made possible by interactive \gls{vr}.\\

\noindent Background knowledge required to grasp the various aspects of this study will be outlined before a selection and comparison of existing literature on, or related to, the subject in Chapter \ref{chap::sota-2}.\\
The subject of the study, and proposed solution theretowards will be detailed in Chapter \ref{chap::design-3}, while the technical contributions made throughout this project are showcased in Chapter \ref{chap::implementation-4}.\\
The proceedings of the study and its results will be looked into in Chapter \ref{chap::evaluation-5}, before closing this dissertation with a critical look at the study and its results, and suggestions for directions research building upon this paper might be worth taking in Chapter \ref{chap::conclusion-6}.
